% In C. Verrem Actio Secunda I (in-verrem-2-1) --- English translation
% Translated by Claude (Anthropic) from the Latin (Perseus, Clark text).
% Style guide: STYLE.md.
%
% This is the first of the five published books of the Actio Secunda,
% the written prosecution Cicero published after Verres withdrew into
% exile rather than face the second hearing.

\ciceroSection{1} I think none of you, gentlemen, is unaware that during these days the talk of the multitude and the opinion of the Roman people has been that Gaius Verres would not answer in the second hearing and would not appear at the trial. Which rumour had spread not only because he had certainly resolved and made up his mind not to appear, but also because no one supposed any man so audacious, so senseless, so shameless, that, having been convicted of such unspeakable charges by so many witnesses, he would dare to look the jurors in the face or show his own face to the Roman people.

\ciceroSection{2} It is the same Verres he always was --- as forward to dare, so prepared to listen. He is here, he answers, he is defended; he does not even leave himself this much, that, when he is being held caught manifestly in the most disgraceful matters, if he were silent and absent, he might at least seem to have sought a modest end to his shamelessness. I bear it, gentlemen, and I do not take it ill that I shall reap the fruit of my labour, and you the fruit of your virtue. For if he had done what he had previously resolved --- not to be present --- it would have been somewhat less known than I needed it to be how I had laboured in preparing and constructing this prosecution; and your praise, gentlemen, would have been thin indeed and obscure.

\ciceroSection{3} Nor does the Roman people expect this from you, nor can it be content with this --- that the man who refused to be present should be condemned, and that you should be brave only against one whom no one dared to defend. On the contrary: let him be present, let him answer; let him be defended with the highest resources, with the highest zeal of the most powerful men. Let my diligence contend with the greed of all of them; your integrity with his money; the constancy of the witnesses with the threats and power of his patrons. Then at last all those things will appear to be conquered, when they have come into contention and rivalry. If he had been condemned in absence, he would seem to have looked not so much to his own advantage as to grudging your praise.

\ciceroSection{4} For no greater safety for the commonwealth can be found at this moment than that the Roman people understand that --- the prosecutor having diligently rejected the jurors --- the allies, the laws, the commonwealth can best be defended by senatorial council. Nor can any greater ruin come to all our fortunes than that, by the opinion of the Roman people, the standard of truth, integrity, faith, religion be judged absent from this order.

\ciceroSection{5} So it seems to me, gentlemen, that I have undertaken a great and most ailing and almost given-up part of the commonwealth, and that I have served not so much my own praise and reputation in this as yours. For I have come to the lifting of the resentment against the courts and the removal of their reproach: that, when this case has been judged according to the will of the Roman people, in some part the authority of the courts may seem to be re-established by my diligence; while if it should be wrongly judged, an end may at last be set to the controversy over the courts.

\ciceroSection{6} For without doubt, gentlemen, in this case that very thing is brought into crisis. The defendant is most guilty. If he is condemned, men will cease to say that money has greatest power in these courts. If he is acquitted, we shall cease to refuse the transfer of the courts. Although about his acquittal he himself no longer hopes, nor does the Roman people fear, there are some who wonder at his singular shamelessness --- that he is here, that he answers. To me, considering the rest of his audacity and madness, this does not seem wonderful. For he has committed many things impiously and unspeakably against gods and men, by the punishments of which crimes he is harried and led from his right mind and judgement.

\ciceroSection{7} The penalties of those Roman citizens drive him headlong --- whom he in part struck with the axe, in part killed in chains, in part hung upon the cross while they invoked the laws of liberty and citizenship. The household gods drag him to punishment, since he was found who would tear sons out of their parents' embrace and lead them to death, and who would demand from the parents the price of the burial of their children. The religions and ceremonies of all the rites and shrines violated, the images of the gods --- which not only have been carried off from their temples but lie thrust away in darkness and hidden by him --- do not allow his mind to stand still without madness and frenzy.

\ciceroSection{8} Nor does this man seem to me to be offering himself only to condemnation; nor is he content with this common penalty for greed --- a man who has bound himself by so many crimes. His monstrous and intolerable nature requires some singular punishment. It is not only sought, that with his condemnation his goods be restored to those from whom they were torn, but the religion of the immortal gods must be expiated, and the torment of Roman citizens and the blood of many innocents must be paid for by his punishment.

\ciceroSection{9} For we have brought into your court not a thief but a robber, not an adulterer but a stormer of chastity, not a sacrilegious man but an enemy of all sacred things and religion, not an assassin but the most cruel butcher of citizens and allies. So I think this is the one defendant within human memory whom it would be advantageous to have condemned. For who does not understand that, even if he is acquitted against the will of gods and men, he can in no way be torn out of the hands of the Roman people? Who does not see clearly that we shall be reckoned to have come off splendidly if the Roman people shall be content with this man's punishment alone, and shall not decide thus: that this man has not committed a greater crime against himself --- when he plundered shrines, when he killed so many innocent men, when he afflicted Roman citizens with death, torture, and the cross, when he released the leaders of the pirates after taking money --- than those, if any, who, having sworn, by their vote shall have freed this man, covered with so many great unspeakable crimes?

\ciceroSection{10} There is not, there is not in this man any room for any erring, gentlemen. He is not the kind of defendant, this is not the moment, this is not the council --- I fear lest I seem to speak too arrogantly before such men --- this is not even the prosecutor through whom a defendant so guilty, so ruined, so convicted, can either be secretly slipped away or torn from us with impunity. Shall I not prove to these jurors that Gaius Verres took money against the laws? Will such men endure that they did not credit so many senators, so many Roman knights, so many cities, so many honourable men from so distinguished a province, the letters of so many peoples and private men --- that they resisted so great a will of the Roman people?

\ciceroSection{11} Let them endure it: we shall find, if we can bring this man alive to another trial, those before whom we may prove that this man in his quaestorship turned aside the public money given to Gnaeus Carbo the consul; we shall find those whom we may persuade that this man, in another's name --- as you learned in the first hearing --- took money from the urban quaestors. There will be those who will censure his audacity in this too, that under several names he subtracted as much as suited him from the head of the tithe-grain. There will perhaps be those, gentlemen, who will think that embezzlement of his most keenly to be punished, in that he did not hesitate to carry off the monuments of Marcus Marcellus and Publius Africanus --- which in their name, but in fact were and were considered the Roman people's --- from the most religious shrines and from the cities of allies and friends.

\ciceroSection{12} Suppose he has come out from the embezzlement trial too: let him meditate about the pirate captains whom he released after taking money, let him see what to answer about those whom he kept hidden in his house in their place; let him consider not only how he can heal our charge but how to deal with his own confession. Let him remember that, in the previous hearing, stirred by the angry and hostile shout of the Roman people, he confessed that the pirate leaders had not been struck with the axe; that he was already then afraid lest it be brought against him as a charge that he had freed them with money. Let him admit --- which cannot be denied --- that, being a private man, he kept the pirate captains alive and unharmed in his house at Rome, until I made him stop. If he can show in that treason trial that this was permitted to him, I will admit it ought to have been allowed. Suppose he escapes from this too: I will go where the Roman people has long been calling me.

\ciceroSection{13} For of the right of liberty and of citizenship the Roman people thinks the judgement is its own --- and rightly so. Let him by all means break through the senatorial councils with his violence; let him burst the courts of all sorts; let him fly out from your strictness: believe me, he will be held by tighter snares before the Roman people. The Roman people will believe these Roman knights, who, produced before you as witnesses earlier, said in their own sight that a Roman citizen, who offered honourable men as bondsmen, was hung on the cross by him.

\ciceroSection{14} The thirty-five tribes will believe that most weighty and most distinguished man, Marcus Annius, who said that in his presence a Roman citizen was struck with the axe. The Roman people will hear that leading man, the Roman knight Lucius Flavius, who in his testimony said that his close friend Herennius, a businessman from Africa, when more than a hundred Roman citizens at Syracuse recognised him and defended him with tears, was struck with the axe. Lucius Suettius will give his good faith and authority and religion --- a man endowed with all distinctions, who said under oath before you that many Roman citizens were most cruelly put to death by violence, by his order, in the stone quarries. When I conduct this case from a higher position, by the kindness of the Roman people, I do not fear that any force can snatch this man from the votes of the Roman people, or that any service of my aedileship can be more ample or more pleasing to the Roman people.

\ciceroSection{15} Wherefore let all attempt all things in this trial. There is now nothing left in this case in which any man can err, gentlemen, except at your own peril. My method has been recognised in past matters; in what remains, it has been tested and provided for. I showed my zeal for the commonwealth at the time when, after a long interval, I brought back the old custom; and at the request of allies and friends of the Roman people --- who are also my close associates --- I laid the indictment of a most audacious man. Which deed of mine the most chosen and most distinguished men --- of whom several were among you --- so approved that to the man who had been Verres's quaestor, and was wronged by Verres and was pursuing just enmities, they refused not only the laying of the indictment, but even the subscription, when he asked for it.

\ciceroSection{16} I went to Sicily for the inquiry. In which business my industry was shown by the speed of my return, my diligence by the multitude of documents and witnesses, and my modesty and conscience by the fact that, when I had come as a senator to the allies of the Roman people --- I, who had been quaestor in that province --- I lodged with my own friends and connections rather than with those who had asked me for help, the defender of the common cause. To no one was my coming a labour or expense, either publicly or privately. I used in inquiry only as much force as the law gave me, not as much as the zeal of those whom this man had harassed could have given me.

\ciceroSection{17} When I returned to Rome from Sicily --- though Verres and his friends, refined and city-bred men, had spread rumours of this kind, that the witnesses might be slowed up: that I had been led off from a real prosecution by a great sum of money --- although it was credited by no one (because both from Sicily there were witnesses who had known me as quaestor in the province, and from here the most distinguished men, who as they themselves are well known, so know each one of us best) --- yet I feared lest anyone should doubt my fidelity and integrity, until we came to the rejection of the jurors. I knew that, in rejecting jurors in our own memory, some had not avoided the suspicion of a deal, while in the prosecution itself their industry and diligence was approved.

\ciceroSection{18} I rejected jurors in such a way that this is agreed: that since the present state of the commonwealth which we now use, no council of similar splendour and dignity has existed. Which praise this man says is common between him and me. He, when he had rejected Publius Galba as juror, kept Marcus Lucretius; and when his patron asked him why he had let his closest friends, Sextus Peducaeus, Quintus Considius, Quintus Junius, be rejected, he replied: "Because I knew them in giving judgement to be too independent in their voting and view."

\ciceroSection{19} So with the jurors rejected, I now hoped that my burden was shared with you. I thought that not only to those who knew me but to those who did not, my fidelity and diligence had been proved. In which I was not deceived. For at my elections, when this man used unlimited bribery against me, the Roman people judged that the man's money, which against my faith had had no power with me, ought to have no power with them against my honour. On which day, gentlemen, when you first sat as called against this defendant, who was so unjust to this order, who was so eager for new things, for new courts and new jurors, that he was not moved by your appearance and your assembling?

\ciceroSection{20} When in this matter your dignity rewarded me for my diligence, this is what I gained: that in the one hour in which I began to speak, I cut off from a defendant audacious, moneyed, lavish, and ruined, the hope of corrupting the trial; that on the first day, with so great a number of witnesses cited, the Roman people judged that, with this man acquitted, the commonwealth could not stand; that the second day took from this man's friends and defenders not only the hope of victory but the will to defend; that the third day so cast the man down that, by feigning illness, he debated not what he should answer but how not to answer. Then on the remaining days, by these charges, by these witnesses both city and provincial, he was so overwhelmed and crushed that, with the days of the games interposed, no one judged him postponed for further hearing, but condemned.

\ciceroSection{21} Wherefore, so far as concerns me, gentlemen, I have won. For I have not coveted the spoils of Gaius Verres, but the reputation of the Roman people. It was for me to come with a cause to the prosecution: what cause was more honourable than to be appointed and chosen as defender by so distinguished a province? To take counsel for the commonwealth: what so suits the commonwealth as in such great resentment against the courts to bring forward a man by whose condemnation the whole order could be in praise and favour with the Roman people? To show and persuade that a guilty man has been brought up: who in the Roman people is there who has not taken away from the previous hearing this conviction --- that all the crimes, thefts, and disgraces of those previously condemned, if they were brought together into one place, would scarcely be matched and compared with a small part of his?

\ciceroSection{22} Look you to and consult, gentlemen, what concerns your reputation, your estimation, and the common safety. Your splendour ensures that you cannot err without the highest detriment and danger to the commonwealth. For the Roman people cannot hope that there are others in the senate who can rightly judge, if you cannot. It is necessary, when it has despaired of the whole order, to seek another kind of men and another method of giving judgements. If this seems lighter to you because you think it is a heavy and inconvenient burden to give judgement, you ought to understand first that it makes a difference whether you yourselves throw off this burden, or whether, because you could not prove your fidelity and conscience to the Roman people, the power of giving judgement is for that reason taken from you. And then to consider also: in what danger we shall come before those jurors whom, on account of hatred of us, the Roman people shall have wished to give judgement on us.

\ciceroSection{23} I will tell you what I have understood, gentlemen. Know that there are some men whom such a hatred of our order possesses that they openly say that they wish this man --- whom they know to be most dishonest --- to be acquitted on this one count alone: that the courts may, with disgrace and shame, be taken from the senate. This has compelled me to deal with you in many words, gentlemen --- not my fear about your fidelity, but a new hope of theirs which, having suddenly drawn Verres back from the gate to the trial, made some suspect that not without cause his plan had been so suddenly changed.

\ciceroSection{24} Now, lest Hortensius use a new kind of complaint and say that a defendant is being crushed about whom the prosecutor says nothing; that nothing is so dangerous to the fortunes of the innocent as for adversaries to be silent; and that he should not praise my ability differently than I wish, when he says that, if I had said many things, I would have lifted up the man against whom I was speaking, but that, because I did not say them, I had ruined him: I shall give way to him. I will use a continuous oration --- not because this is now necessary, but that I may test whether he bears it more troublesomely that I was silent then or that I am speaking now.

\ciceroSection{25} Here you will perhaps be diligent to see that I do not remit any hour from my legal hours. Unless I shall have used up every moment which the law has granted me, you will complain, you will invoke the faith of gods and men: that Gaius Verres is being circumvented because the prosecutor refuses to speak as long as he is allowed. Shall I not be allowed not to use what the law has given me for my own sake? For the time of accusation has been given me for my own sake, that I may be able to set forth in my speech the charges and the case. If I do not use it, I do you no wrong, but I take something from my own right and convenience. "But the case must be heard." For this reason indeed: because otherwise the defendant, however guilty, cannot be condemned. Did you, then, take it ill that something was done by me to keep this man from being condemned? For with the case heard many can be acquitted; with it not heard, no one can be condemned.

\ciceroSection{26} "I take away the second hearing." That which the law has in itself most troublesome, that the case should be pleaded twice --- which is established more for my sake than yours, or no more for yours than mine. For if to plead twice is convenient, surely it is common to both. If it is needful that he who spoke later be refuted, it has been established for the prosecutor's sake that the matter should be conducted twice. But, I take it, Glaucia first proposed that the defendant be put off for a second hearing. Before that, judgement could be given at the first or "amplius" pronounced. Which then do you think the milder law? I think the older one, by which one could be either quickly acquitted or slowly condemned. I am restoring to you that lex Acilia, by which many were once accused, the case once pleaded, the witnesses once heard, and condemned --- by no means with charges so manifest or so great as those by which you are convicted. Suppose you were pleading your case not under this so atrocious law, but under that mildest one. I shall accuse; you will answer; with the witnesses produced, I shall send the matter to the council in such a way that, even if the law gave the power of "amplius," yet they would think it disgraceful for them not to give judgement at the first.

\ciceroSection{27} But if it is necessary for the case to be heard --- has it not been heard enough? We pretend, Hortensius, what we have often experienced in speaking. Who paid great attention to us in this kind of case, where something is said to have been snatched or carried off by anyone? Is not all the expectation of the jurors either in the documents or in the witnesses? I said in the first hearing that I would make plain that Gaius Verres had taken from Sicily forty million sesterces against the laws. What? Should I have made this plainer if I had narrated it thus? "There was a certain Dio of Halaesa, who, when his son was an heir of his kinsman to a very great inheritance under the praetor Gaius Sacerdos, had at that time no business or controversy. Verres, as soon as he touched the province, immediately gave letters at Messana, summoned Dio, set against him slanderers from his own bosom who said that the inheritance had fallen to Erucine Venus; about this matter he showed he himself would judge."

\ciceroSection{28} I can next unfold the whole matter, then say in conclusion what happened: that Dio counted out a million sesterces to win the most certain case; that besides, his herds of mares were carried off by him, that the silver, the embroidered tapestries he had he saw to be carried off. These things, neither when I said them nor when you denied them, would our speech be of great moment. At what time, then, would the juror prick up his ear and apply his mind? When Dio himself came forward; when others who were then in Sicily and present at Dio's affairs; when through those very days when Dio was pleading the case, it was found he had taken money on loan, called in his debts, sold properties; when documents of good men were produced; when those who gave money to Dio said that they had then heard those funds were being taken to be given to Verres; when Dio's friends, hosts, patrons, the most honourable men, said they had heard the same things.

\ciceroSection{29} I take it, when these things were happening, then you would listen --- as you have listened --- then the case would seem truly to be conducted. Thus I have done all things in the previous hearing in such a way that on no charge was there any in which any of you needed a continuous prosecution. I deny that anything has been said by the witnesses which was either obscure to any of you or required any orator's eloquence. For you remember that I myself acted in such a way that, in interrogating witnesses, I set forth and unfolded all charges; so that, when I had placed the whole matter in plain view, I then at last interrogated the witness. So not only do you, who must judge, hold our charges; but the Roman people too has learned the whole prosecution and the case. Although I speak of my deed as if I had done that of my own will rather than driven by your wrong.

\ciceroSection{30} You interposed a prosecutor who, when I had asked for myself only one hundred and ten days for Sicily, asked one hundred and eight for himself for Achaia. When you had snatched from me three months most fitted for action, you thought that I would yield up to you all the rest of this year's time --- so that, when I had used up our hours, you, with two sets of games interposed, would answer on the fortieth day after; then time would be drawn out so that we would come from Manius Glabrio the praetor and from a great part of these jurors to another praetor and other jurors.

\ciceroSection{31} If I had not seen this --- if all my acquaintances and strangers had not warned me that this was being done, this was being plotted, this was being laboured at --- that the matter should be put off to that moment --- I might, I suppose, fear that, if I wished to use my hours in prosecuting, my charges would not last out, my speech would fail, my voice and strength would give way; that I should not be able to prosecute twice the man whom no one had dared to defend in the first hearing. I have proved my plan both to the jurors and to the Roman people. There is no one who thinks the wrong and shamelessness of these men could have been opposed in any other way. For with what folly should I have been, if --- when those who had bought to rescue Verres set down the day in their bond ("if the matter goes to the council after the first of January") --- I should have fallen on that day myself, when I could have avoided it?

\ciceroSection{32} Now I must take careful account of the time given me for speaking, since it is in my mind to set out the whole case. So I shall pass over the first act of his most disgraceful and flagitious life. He shall hear nothing from me about the disgraces of his boyhood, nothing from that impure youth of his --- which kind it was, you either remember, or, from the man whom he has produced as most like himself, you can recognise. I will pass over everything that seems shameful to me to mention; nor will I consider only what it suits him to hear, but what it suits me to say. I beg you to grant this and to allow my modesty: that I may keep silent on some part of his shamelessness.

\ciceroSection{33} All that time which was before he came to magistracies and the commonwealth, let it be free and unbound for me. Let his nightly orgies and vigils be passed over in silence; let no mention be made of pimps, gamblers, panders. Let the losses, the disgraces which his father's estate, his own youth bore, be passed by. Let him profit by the silence about his earlier infamy; let his rest of his life suffer me to make this great loss of charges.

\ciceroSection{34} You were quaestor to Gnaeus Papirius the consul, fourteen years ago. From that day to this, what you have done, I summon to court: not an hour will be found empty of theft, crime, cruelty, disgrace. These are the years spent in the quaestorship and the legateship of Asia and the city praetorship and the Sicilian praetorship. So this same fourfold division will be of my whole prosecution. As quaestor, you drew the province by senatorial decree: the consular fell to you, that you should be with the consul Gnaeus Carbo and hold that province. There was at that time dissension among the citizens, about which I will say nothing of what you ought to have felt. This one thing I say: at such a time and in such a lot, you ought to have decided which side you preferred to feel for and to defend. Carbo bore it ill that a man of singular luxury and idleness had fallen to him as quaestor; nevertheless he honoured him with all kindnesses and offices. To make this short: the money was assigned and counted out. The quaestor set out for the province; he came to the consular army in Gaul as expected, with money. As soon as the first opportunity seemed to be given him --- learn how a man begins his magistracies and his administration of the commonwealth --- the public money turned aside, the quaestor deserted his consul, his army, his lot, and his province.

\ciceroSection{35} I see what I have done. He pricks up his ears; he hopes that some breath can be blown his way in this matter, by the will and defence of those to whom the name of the dead Gnaeus Carbo is hateful, with whom he hopes that abandonment and betrayal of his consul will be welcome. As if he had done this from a desire to defend the nobility or zeal for that party, and not most openly plundered the consul, the army, and the province, and run away because of the most shameless theft! For it is obscure and his deed of such a kind that anyone might suspect that Gaius Verres, because he could not endure new men, had crossed over to the nobility, that is, to his own kind --- and had done nothing for the sake of money!

\ciceroSection{36} Let us look at his accounts, how he rendered them: now he himself will show why he left Gnaeus Carbo, now he will betray himself. First learn the brevity. "Received," he says, "two million two hundred thirty-five thousand four hundred seventeen sesterces. Paid out for pay, grain, legates, the deputy quaestor, the praetorian cohort, one million six hundred thirty-five thousand four hundred seventeen sesterces. Remaining: at Ariminum, six hundred thousand sesterces." Is this how to render accounts? Have I or you, Hortensius, or anyone in the world, ever rendered them in this fashion? What is this? What shamelessness, what audacity? What example, out of so many men's accounts rendered, is of this kind? But that six hundred thousand sesterces, which he could not falsely set down to whom they had been given, which he writes were left at Ariminum, those very six hundred thousand which were left over --- Carbo did not touch them, nor did Sulla see them, nor were they brought back to the treasury. He chose Ariminum as his town --- a town which, when he was rendering his accounts, had been crushed and plundered. He did not suspect, what he will now perceive, that we have witnesses enough left from that calamity of the Ariminenses for this matter. Read on.

\ciceroSection{37} "To Publius Lentulus and Lucius Triarius, urban quaestors, the matter of the accounts was rendered." Read: "By decree of the senate." So that he might be permitted to render his account in this fashion, he became a Sullan suddenly --- not that honour and dignity might be restored to the nobility. But if you had fled empty from there, still that wicked desertion and unholy betrayal of your consul would be judged a crime. "Carbo was a bad citizen, a wicked consul, a seditious man." So he was to others; when did he begin to be so to you? After he had committed to you his money, his grain supply, all his accounts, and his army. For if he had earlier displeased you, you would have done what Marcus Piso did the next year. When the quaestorship had fallen to him to Lucius Scipio the consul, he did not touch the money, did not set out for the army; what he felt about the commonwealth, he felt without injuring his own faith, the custom of our ancestors, or the bond of the lot.

\ciceroSection{38} For if we wish to confound and confuse all these things, we shall make all life dangerous, hateful, and hostile --- if the lot will have no religion, the bond of mutual fortune (whether favourable or doubtful) no companionship, the customs and institutions of our ancestors no authority. He is the common enemy of all who has been the enemy of his own. No wise man has ever thought a betrayer should be trusted. Sulla himself, to whom this man's coming ought to have been most welcome, removed him from himself and from his army; he ordered him to be at Beneventum, with those whom he understood to be most friendly to his party, where he could in nothing harm the highest matter and cause. He afterwards bestowed rewards on him generously, granted him certain goods of the proscribed in the Beneventan country to be plundered. He gave honour as to a betrayer, not faith as to a friend.

\ciceroSection{39} Now however much there are men who hate the dead Gnaeus Carbo, they ought to think not what they wished to happen to him, but what is to be feared by themselves in such a case. This is a common evil, a common fear, a common danger. There are no more hidden snares than those which lie hidden in the pretence of duty or in some name of close association. For one who is openly an adversary you can easily avoid by being on guard; but this hidden, internal, household evil not only does not appear, but even crushes you before you can see and explore it. Is it really so?

\ciceroSection{40} You, when you have been sent as quaestor to the army --- the guardian not only of the money but also of the consul, partner in all matters and counsels, held in the place of a son, as the custom of our ancestors used --- shall you suddenly leave him, desert him, cross over to the adversaries? O crime! O monster to be exported to the farthest lands! For that nature which has committed so great a deed cannot be content with this one crime. It must always plot something of the kind, must always be engaged in similar audacity and treachery.

\ciceroSection{41} So this same man, whom Gnaeus Dolabella afterwards, when Gaius Malleolus had been killed, had as deputy quaestor --- I do not know whether this connection was not even greater than that with Carbo, and the judgement of the will ought not to count more than that of the lot --- the same man was to Gnaeus Dolabella as he had been to Gnaeus Carbo. For he turned the charges that had weight against himself onto Dolabella, and brought the whole case of his patron over to Dolabella's enemies and prosecutors; he himself spoke a most hostile and most dishonest testimony against the man to whom he had been legate, to whom he had been deputy quaestor. That wretched Gnaeus Dolabella was burnt down --- both by his unspeakable betrayal and by his dishonest and false testimony --- and most of all by the resentment against this man's own thefts and disgraces.

\ciceroSection{42} What shall you do with this man? For what hope shall you reserve so treacherous, so intolerable a beast --- who in Gnaeus Carbo neglected and violated the lot, in Gnaeus Dolabella the will, and both not only deserted but betrayed and assaulted? Do not, gentlemen, please, weigh these charges by the brevity of my speech rather than by the magnitude of the matters themselves; for I must necessarily hurry to set out for you everything that has been laid down for me.

\ciceroSection{43} Wherefore, his quaestorship being demonstrated and the theft and crime of his first magistracy seen through, attend to the rest. In which I shall pass over that period of the Sullan proscriptions and plunderings; nor will I let him take any defence for himself out of the common calamity. I shall accuse him of his own definite and proper crimes. Wherefore, with all this Sullan period excluded from the prosecution, learn his splendid legateship.

\ciceroSection{44} After Cilicia was decreed as Gnaeus Dolabella's province --- O immortal gods, with what greed, by what intrigues, did this man storm that legateship for himself! That was the beginning of Gnaeus Dolabella's greatest calamity. For as he set out, wherever he made his journey, he was such that he seemed to spread not as a legate of the Roman people but as a kind of calamity. In Achaia --- I will pass over all smaller things, of which something similar perhaps another man too has done at some time. I will say nothing but what is singular, what, if it were said against another defendant, would seem incredible --- he demanded money of the magistrate of Sicyon. Let this not be a charge against Verres: others have done the same. When the man would not give it, he punished him --- dishonestly, but not unheard of.

\ciceroSection{45} See what kind of punishment --- you will ask from what kind of man you judge him. He ordered a fire to be made of green and damp wood in a narrow place; there he left a free-born man, noble in his own city, an ally and friend of the Roman people, half-dead, tormented by the smoke. As for the statues, the painted panels he carried off from Achaia --- I will not speak in this place; I have another place set apart for showing his greed for these things. At Athens, you have heard that a great mass of gold was carried off from the temple of Minerva; this was said in Gnaeus Dolabella's trial. Said? Even assessed. You will find that Gaius Verres was not a sharer in this counsel, but the leader.

\ciceroSection{46} He came to Delos. There from the most religious shrine of Apollo by night he secretly took away the most beautiful and most ancient statues, and saw to it that they were thrown into his cargo ship. The next day, when those who lived at Delos saw the shrine plundered, they took it ill; for so great is the religion and antiquity of that shrine among them that they think Apollo himself was born in that very place. Yet they did not dare say a word, lest the matter perhaps concern Dolabella himself. Then suddenly the greatest tempests arose, gentlemen, so that not only could Dolabella, although he wished, not set out, but he could scarcely stay in the town: so great were the waves cast up. Here that ship of this brigand, loaded with sacred statues, was driven out and broken by the waves; on the shore those statues of Apollo are found. By Dolabella's order they are restored. The tempest is calmed; Dolabella sets out from Delos.

\ciceroSection{47} I do not doubt that, although there has never been any feeling of humanity in you, never any thought of religion, yet now in your fear and danger your own crimes come to your mind. Can any hope of safety be conveniently shown you, when you remember how impious, how wicked, how unspeakable you have been towards the immortal gods? Did you dare to plunder Apollo of Delos? Did you dare to lay impious and sacrilegious hands on that temple, so ancient, so sacred, so religious? If in your boyhood you were not so trained in those arts and disciplines that you should learn and know what has been entrusted to letters --- not even afterwards, when you came to those very places, were you able to learn what is handed down by memory and writing:

\ciceroSection{48} that Latona, after long wandering and flight, when pregnant and her time of delivery now upon her, fled to Delos and there bore Apollo and Diana? From which opinion of men that island is thought sacred to those gods, and the authority of its religion is so great, and always has been, that not even the Persians, when they had declared war on all of Greece, gods and men, and brought a fleet of a thousand ships to Delos, attempted to violate or touch anything. Did you, most dishonest and most senseless of men, dare to plunder this shrine? Was there ever a greed so great that it could overcome so great a religion? And if you did not think of these things then, do you not now even remember that no evil is so great but it has long been owed you for your crimes?

\ciceroSection{49} When he came to Asia --- why should I speak of his receptions, his lunches, dinners, horses, and gifts? I will not deal with Verres on his daily crimes. At Chios, I say, he carried off the most beautiful statues by force; likewise at Erythrae and Halicarnassus. At Tenedos --- I pass over the money he plundered --- Tenes himself, who is held by the Tenedians the most sacred god, who is said to have founded their city, from whose name Tenedos is named --- this very Tenes, I say, most beautifully wrought, whom you once saw in the Comitium, he carried off, with the great groan of the city.

\ciceroSection{50} As for that storming of the most ancient and most noble shrine of Samian Juno --- how grievous to the Samians, how bitter to all Asia, how known to all, how unheard of by no one of you! On which storming, when ambassadors had come from Samos to Gaius Nero in Asia, they took back this answer: that complaints of this kind, which concerned legates of the Roman people, ought to be referred not to the praetor but to Rome. What documents did this man take from there, what statues! Which I myself recently recognised in his house, when I had come for the sake of sealing them up.

\ciceroSection{51} Where now, Verres, are these statues? Those, I mean, which we recently saw at your house, set out by all the columns, in all the spaces between the columns, and finally in your park, under the open sky. Why, as long as you thought another praetor with the jurors you had drawn for the substitution would be entering on the council about you --- so long they were at home? After you saw that we preferred to use our witnesses than your hours, you left no statue at home except two which are in the middle of the house, which were carried off from Samos itself. Did you not think I would summon as witnesses on this matter your closest friends, who had always been at your house, of whom I might inquire whether they knew the statues had been there which were not now there?

\ciceroSection{52} What did you suppose those men would judge of you, when they saw that you were now fighting not against your prosecutor but against an inheritance-officer and a confiscation-broker? Concerning which matter, you heard Charidemus of Chios speak in evidence in the previous hearing: that, when he was a trierarch and was escorting Verres from Asia by Dolabella's order, he was at Samos with him, and that he then knew the shrine of Juno and the town of Samos were plundered. And afterwards he pleaded his case publicly before his fellow citizens at Chios, the Samians being his accusers. He was acquitted, because he made it plain that what the legates of the Samos said pertained to Verres, not to him.

\ciceroSection{53} Aspendus --- you know it is an old and noble town in Pamphylia --- is full of the best statues. I will not say this statue was taken from there and that one. I say this: that you left no statue at Aspendus, Verres. Everything from the shrines, from public places, openly, in everyone's sight, was carried out and exported on wagons. And even that famous Aspendian harpist, of whom you have often heard --- the one whom the Greeks proverbially say "plays everything within" --- he carried off and set up in his innermost chambers, so that he might seem to have surpassed even him in his own art.

\ciceroSection{54} At Perga we know there is the most ancient and most sacred shrine of Diana. That too has been stripped and plundered by you; from Diana herself, what gold she had, has been pulled off and carried away, I say. What, in evil hour, is this so great audacity and madness? For the cities of allies and friends to which you came under the right and name of legate --- if you had invaded them by force, with army and command, yet, I take it, the statues and ornaments which you carried off from those cities you would have carried off not to your own house, nor to the country-houses of your friends, but to Rome, to the public.

\ciceroSection{55} What shall I say of Marcus Marcellus, who took Syracuse, that most splendid of cities? What of Lucius Scipio, who waged war in Asia and conquered Antiochus, the most powerful king? What of Flamininus, who subdued King Philip and Macedon? What of Lucius Paulus, who overcame King Perses by force and valour? What of Lucius Mummius, who took the most beautiful and most ornate city of Corinth, full of every kind of treasure, and joined many cities of Achaia and Boeotia under the rule and dominion of the Roman people? Whose houses, when they flourished in honour and virtue, were empty of statues and panel paintings; but the whole city and the temples of the gods and all parts of Italy we see adorned with their gifts and monuments.

\ciceroSection{56} I am afraid lest these things may seem to someone perhaps too ancient and outworn now; for then all were equally of this kind, so that this praise of singular virtue and innocence seems to be of those men, indeed of those times, not of these. Publius Servilius, a most distinguished man, with the greatest deeds done, is here to give his vote on you. He took Olympus by force, by troops, by counsel, by valour --- an ancient city, increased and adorned with all things. I produce a recent example of a most brave man. For Servilius, commander of the Roman people, took Olympus, an enemy city, after you in those same parts as a quaestorian legate had pacified towns of allies and friends seen to be plundered and harassed.

\ciceroSection{57} What you carried off by crime and brigandage from the most religious shrines, we cannot see except in your own house and those of your friends. The statues and ornaments which Publius Servilius carried off from an enemy city taken by force and valour, by the law of war and the right of imperator, he brought to the Roman people, conveyed in his triumph, saw to be inscribed on the public tablets at the treasury. Learn from the public letters the diligence of so distinguished a man. Read: "Accounts rendered by Publius Servilius." You see that not only the number of the statues but the size, shape, and posture of each one is set down in the documents. Surely the joy of valour and victory is greater than this pleasure which is taken from lust and greed. I say Servilius keeps the booty of the Roman people much more carefully marked and recorded than you keep your thefts.

\ciceroSection{58} You will say that your statues and panel paintings were also an ornament to the city and the Forum of the Roman people. I remember; with the Roman people I saw the Forum and the Comitium adorned for show with magnificent ornament, but bitter and mournful to the senses and the mind. I saw everything ablaze with your thefts, the booty of the provinces, the spoils of allies and friends. At which time, gentlemen, this man got the greatest hope for his other crimes. For he saw that those who wished to be called the masters of the courts were the slaves of these greeds.

\ciceroSection{59} The allies and foreign nations then for the first time threw away all hope of their affairs and fortunes, because by chance there were at Rome at that time very many ambassadors from Asia and Achaia, who venerated in the Forum the statues of the gods carried off from their own shrines, and likewise, when they recognised other statues and ornaments, looked at them in tears, one in this place, another in that. The talk of all these we then heard was that there was no reason for anyone to doubt about the destruction of the allies and friends, when indeed they saw in the Forum of the Roman people, in that place where formerly those who had done injuries to allies were accused and condemned, there openly placed those things which had been taken and torn from the allies by crime.

\ciceroSection{60} Here I do not think he will deny that he has very many statues, innumerable panel paintings; but, I take it, he is accustomed sometimes to say that what he has seized and stolen he bought, since indeed into Achaia, Asia, Pamphylia, at public expense and in the name of an embassy, he was sent as a buyer of statues and panel paintings. I have received all the documents both of him and of his father, which I have read most diligently and digested: his father's, while he lived; yours, as long as you say you kept them. For in him, gentlemen, you will find this novel thing. We hear that some never kept books --- which is the false opinion men have of Antonius, for he kept them most diligently; but let it be a certain kind of conduct, by no means to be approved. We hear that another did not keep them from the start, but kept them from some time onward; this also has some method. But this is novel and ridiculous: that this man answered us when we asked him for his books, that he had kept them up to the consulship of Marcus Terentius and Gaius Cassius, but afterwards had stopped.

